% Manuscript rectification map: all load-bearing locations that must
% carry the two-stage factorisation + universal positive-geometry grammar.
% Not a manuscript chapter --- an internal inventory.
% Author: Raeez Lorgat.

\providecommand{\kBKM}{\kappa_{\mathrm{BKM}}}
\providecommand{\kch}{\kappa_{\mathrm{ch}}}
\providecommand{\kcat}{\kappa_{\mathrm{cat}}}
\providecommand{\kfib}{\kappa_{\mathrm{fiber}}}

\section*{Manuscript rectification map --- lossless propagation}
\label{wn:sec:manuscript-rectification-map}

\emph{Lossless rectification. Fourteen load-bearing locations must
each carry the two-stage factorisation, the universal positive-geometry
grammar, and the eight-form catalogue. Existing mathematical content
is invariant; nothing is deleted. New material is inserted as
additional framing, with the existing material reinterpreted through
the new grammar. When a prior claim and a new claim describe the
same mathematical content under different framings, both appear:
the prior claim in its established form, the new framing as the
unifying lens. Below is the inventory: file, augmentation target,
and the guarantee that no mathematical content is lost.}

\subsection*{Lossless principle}

Every claim, proof, construction, formula, citation, and example
presently in the manuscript remains in the manuscript. New
material is additive. The new grammar \emph{reorganises the
hierarchy of exposition}; the content at the bottom of every
reorganised hierarchy is preserved verbatim. Where a section is
relabelled, the label is added, the section remains. Where an
interpretation shifts, the new interpretation sits alongside the
old with a bridge sentence locating the shift.

\subsection*{Front matter}

\paragraph{\texttt{abstract}.} Prepend the content of
\texttt{notes/platonic\_abstract.tex} as a new foregrounding
paragraph. Any existing abstract material remains verbatim as
subsequent paragraphs; the platonic abstract gives the unifying
lens, the existing abstract's content specifies particular
programme achievements.

\paragraph{\texttt{chapters/frame/preface.tex}.} Insert at the head
(before the existing opening section but after any metadata
preamble) the content of \texttt{notes/platonic\_preface\_opening.tex}.
The existing preface material remains in its current position,
in full. The insertion adds four new sections that announce the
grammar; the existing cascade table, pentagon diagrams, six-route
catalogue, CHL orbifold content, $\rho^{R_i}$ stratification and
all other present material is retained verbatim. Where the new
material refers to an existing table or diagram, the existing
object is cited by its current label; no renumbering that removes
content takes place. The interpretation bridge (``six applications
of $\Phi_3$'' $\leadsto$ ``six $(\Sigma_2, C)$-specialisations of
one $E_3$-hFA or Stage-1 invariants'') appears as a new remark
inserted near the pentagon diagram; the pentagon diagram itself
remains unchanged.

\paragraph{\texttt{chapters/theory/introduction.tex}.} Insert
\texttt{notes/platonic\_introduction\_opening.tex} as a new Section~1
of the chapter, renumbering the existing sections $1$-through-$n$ to
$2$-through-$(n+1)$. No section is removed, rewritten, or
shortened. The existing introduction material on $A_\infty$-categories,
cyclic pairings, and Hochschild calculus remains verbatim as its
own sections; the platonic introduction gives the unifying lens
under which the foundational material is organised.

\subsection*{Part II: the functor $\Phi$ (lossless)}

\paragraph{\texttt{chapters/theory/cy\_to\_chiral.tex}.} Prepend a
new opening section stating
Theorem~\ref{wn:thm:platonic-intro-stage-1} and
Theorem~\ref{wn:thm:platonic-intro-stage-2}. All existing content
remains in place. The existing proof that
$\Phi(\CoHA(\CC^3)) = Y^+(\widehat{\mathfrak{gl}}_1)$ is
retained verbatim; a bridge remark is added noting that this result
is now a corollary of Stage~1 applied to $\cA = \CoHA(\CC^3)$
composed with Stage~2 at the point-specialisation
$(\mathrm{pt}, \CC)$ --- the two formulations describe the same
mathematical content and both appear. Additionally insert the
worked example
$\Sp_{K3, E}(\cF_{K3 \times E}) \simeq
U^{\mathrm{chir}}(\mathfrak{heis}_{\mathrm{Muk}}) \otimes
U^{\mathrm{chir}}(\fg^{\mathrm{BPS}}_{K3})$ as a new theorem
(building on Wave-16 S1 inscription), adjacent to the existing
$d = 3$ discussion; the existing $d = 3$ discussion is preserved.

\paragraph{\texttt{chapters/theory/quantum\_chiral\_algebras.tex}.}
Insert a new framing paragraph locating the existing
six-dimensional holomorphic Chern--Simons observables construction
within the Stage-1 framework: the observables \emph{are} the
factorisation algebra on $X$, which is the Stage-1 output; the
chiral-on-a-curve shadow is the Stage-2 specialisation of the
same factorisation algebra. The existing chapter body is preserved
verbatim. The new framing paragraph appears as a \textbf{Remark
(two-stage locus)} immediately after the construction statement,
pointing the reader to the preface's Stage-1/Stage-2 decomposition
without altering any existing claim, proof, or formula.

\subsection*{Part III: $E_n$ hierarchy}

\paragraph{\texttt{chapters/theory/e2\_chiral\_algebras.tex}.} Shift
framing: $E_2$-chiral on K3 is the Stage-2 specialisation at
$(\Sigma_1, C) = (S^1 \subset K3, \text{Nakajima cycle})$, not
the native output. The Nakajima--Mukai result becomes a corollary
of the specialisation theorem.

\paragraph{\texttt{chapters/theory/en\_factorization.tex}.}
Stage 1 output is $E_d$-hFA on $X$; Stage 2 specialisation on a
$(d{-}1)$-cycle lands $E_1$-chiral on the curve. Existing material
on $E_n$-operad structure propagates into Stage 1.

\subsection*{Part IV: quantum groups and Drinfeld centre}

\paragraph{\texttt{chapters/theory/quantum\_groups\_foundations.tex}.}
Insert the universal $Y^+(X) =
H^{\bullet}_{\mathrm{eq}}(\Meff(X), \phi_W)$ definition at the
head. The equivariance stratification (toric / reduced / orbifold /
lattice-polarised) is the primary organising principle.

\paragraph{\texttt{chapters/theory/drinfeld\_center.tex}.} The
Drinfeld centre of $\mathrm{Rep}(\Phi_3(\cA))$ is the right adjoint
to the forgetful functor, not an averaging. The centre at $d = 3$
lives on the $E_2$-braided category transported by Stage 2
along $(\Sigma_2, C)$ choice; it is not a native output at $d = 3$.

\subsection*{Part V: the Calabi--Yau landscape}

\paragraph{\texttt{chapters/examples/k3\_chiral\_algebra.tex}.} The
K3 chiral algebra is the Stage-2 specialisation of $\cF_{K3}$ at
$(\Sigma_1, C) = (S^1, \text{Nakajima cycle})$. Add a section on
$Y^+(K3) = H^\bullet(\mathrm{Hilb}(K3), \phi_W) \cong \CoHA(Q_{K3})$
via Kapranov exceptional collection on K3 Mukai quiver.

\paragraph{\texttt{chapters/examples/k3e\_bkm\_chapter.tex}.}
Restructure as four-corner theorem on $K3 \times E$:
\begin{enumerate}[label=(\roman*)]
\item the Pandharipande--Oberdieck reduced Donaldson--Thomas moduli
as positive effective geometry, $\CC^\times_E \times \mathrm{Aut}_s(K3)$-equivariant;
\item the cohomological Hall algebra $Y^+(K3 \times E)$;
\item the BPS Lie algebra $\fg_{\mathrm{BPS}}(K3 \times E)$ via
Davison--Meinhardt;
\item the Borcherds automorphic correction $\Phi_{10} = \Delta_5^2$.
\end{enumerate}
Four distinct constructions, one canonical object.

\paragraph{\texttt{chapters/examples/cy\_d\_kappa\_stratification.tex}.}
Explicit subscripted $\kappa$-accounting across all dimensions.
On $K3 \times E$: $\kcat = 0$ (K\"unneth-multiplicative total
space), $\kch = 0$ (Hodge supertrace), $\kBKM(\Delta_5) = 5$
(Borcherds weight), $\kfib(K3) = 24$ (Mukai lattice rank).
Four values from four distinct constructions; no conflation.
Universal identity $\kBKM(\Phi_N) = c_N(0)/2$ holds across all
eight Gritsenko--Cl\'ery forms with cover selection by weight
integrality.

\paragraph{\texttt{chapters/examples/toric\_cy3\_coha.tex}.} Explicit
conifold and local $\mathbb{P}^2$ positive-half construction via
\v{C}ech descent on quiver-affine charts, with Klebanov--Witten
quiver for the conifold and McKay quiver for $\CC^3/\ZZ_3 \simeq
K_{\mathbb{P}^2}$. Gluing cocycle via cluster mutation.

\paragraph{\texttt{chapters/examples/k3e\_cy3\_programme.tex}.} The
terminal closer of Part V. Assembles the reduced DT positive
geometry, the \v{C}ech / Van den Bergh NCCR substitute
(Serre-equivariant quasi-NCCR on $K3 \times E$, since global NCCR
is obstructed by rank-count and Kummer exceptional collection
failure), and the eight-form catalogue with per-form host
geometry. The six routes to $G(K3 \times E)$ appear here
explicitly re-indexed as: four Stage-1 invariants (Mukai, Hodge,
Igusa-weight, fibre) plus three Stage-2 specialisations (BKM on
$(K3, E)$, Niemeier-twist family, Humbert-monodromy limit).

\subsection*{Part VI: seven faces of $r_{\mathrm{CY}}$}

\paragraph{\texttt{chapters/connections/cy\_holographic\_datum\_master.tex}.}
Reinterpret the seven-face picture via the three-tier stratification:
Tier~(i) Calabi--Yau-datum intrinsics (Mukai, Hodge); Tier~(ii)
Stage-1 invariants of $\cF_X$ (fibre rank);
Tier~(iii) $(\Sigma_2, C)$-specialisations (BKM, Niemeier, Humbert,
CHL-twined). Only Tier~(iii) faces are genuine siblings. The
``algebraic seven'' (bar--cobar, CoHA, coisson, MO, Yangian,
Sklyanin, Gaudin) is an orthogonal slicing of the same object.

\paragraph{\texttt{chapters/connections/geometric\_langlands.tex}.}
The Gaiotto curve for the programme's class-$\cS$ connection is
$\Sigma_{0, 24}$ (24-punctured sphere), not $\Sigma_{2, 0}$
(genus-2 closed): Chacaltana--Distler 2010 Table 3 row 1 gives
$c_{4d} = (5n - 13)/6 = 107/6$ at $n = 24$, $c_{2d} = -214$. The
vertex operator algebra $\cV_{24}$ is the iterated Drinfeld--Sokolov
reduction $H^0_{DS}(L_{-2 + 1/22}(\mathfrak{sl}_2)^{\otimes 22})$,
not $L_{-6}(\mathfrak{e}_8)$: the central charge $-214$ excludes
the $\mathfrak{e}_8$ candidate at $c = -62$.

\paragraph{\texttt{chapters/examples/cy\_c\_beyond\_k3e\_existence\_obstruction.tex}.}
The existence obstruction for $G(X)$ beyond $K3 \times E$ lives
on the geometric axis: global Van~den Bergh non-commutative crepant
resolution fails for compact non-toric Calabi--Yau threefolds by
Seidel--Thomas spherical-twist obstruction (rank-count: the K3
Mukai lattice has rank $24$, no tilting bundle of this rank
exists on K3 by Serre-duality). The substitute is the
Serre-equivariant quasi-NCCR built from tilting objects varying
over the Bridgeland stability manifold with $\bS_X$-intertwining
wall-crossing functors.

\subsection*{Part VII: frontiers}

\paragraph{\texttt{chapters/theory/modular\_trace.tex}.} The
modular trace connects the Stage-2 chiral-algebra character to
the Stage-1 reduced-DT generating function; on $K3 \times E$ at
$(\Sigma_2, C) = (K3, E)$ this is the
Oberdieck--Pixton $-1/\Phi_{10}$ identity.

\subsection*{Cross-referencing and labelling}

Every load-bearing location names the others. The sentence
``the universal positive-geometry grammar of $G(X)$'' appears
in the preface, the introduction, the $K3 \times E$ chapter, and
the Calabi--Yau-to-chiral construction chapter --- each time
referring to the same definition, never redefined. The labels
\texttt{thm:platonic-intro-stage-1},
\texttt{thm:platonic-intro-stage-2},
\texttt{thm:platonic-preface-universal-weight} are the canonical
anchors.

\subsection*{What this rectification is not}

This rectification is not a retrospective commentary on earlier
drafts. The manuscript says what it says; previous versions do
not appear. The mathematical content is invariant under the
reorganisation; the presentation is unified by the grammar.
No theorem, lemma, proposition, corollary, remark, construction,
definition, formula, equation, table, diagram, example, or
citation is removed or shortened.

\subsection*{Operative instruction for any future rectifier}

When editing any of the fourteen locations above: add, never
remove. If a sentence appears to be made redundant by new
material, retain it and add a bridge remark connecting the old
and new framings. If a section title no longer reflects its
content, rename the section and retain the content. If a
theorem is restated under new language, the prior statement
remains as a numbered consequence or remark. Mathematical loss
is forbidden; mathematical augmentation is the task.
